\chapter{CONCLUSÃO}
\label{cap:conclusao}

Este trabalho desenvolveu e avaliou uma abordagem de alocação
de ações brasileiras que integra previsões do XGBoost e
sentimento de manchetes ao modelo Black--Litterman.
A avaliação contemplou as etapas de previsão, formação
das opiniões e distribuição dos pesos, além do desempenho
financeiro das carteiras.

\section{Síntese do Trabalho e Atendimento dos Objetivos}
\label{sec:conclusao_sintese}

O objetivo de implementar uma abordagem híbrida de otimização
de portfólios foi atendido. O sistema reúne dados históricos
de mercado, atributos quantitativos e notícias, transformando
essas informações em expectativas de retorno e decisões
simuladas de investimento. O registro das previsões, das
opiniões e das operações permitiu acompanhar o funcionamento
das etapas e verificar seus resultados.

Na configuração consolidada, o componente quantitativo
utilizou modelos de regressão XGBoost treinados com dados
das vinte ações em conjunto, para horizontes de cinco,
dez e vinte e um pregões. O horizonte de cada ação foi
definido pela acurácia direcional em 2024, considerando
apenas resultados conhecidos antes de 2025. As previsões
foram produzidas pelo procedimento \emph{walk-forward}.

O componente textual classificou as manchetes em positivas,
negativas ou neutras. Esses sinais foram combinados às
previsões quantitativas para formar as opiniões de retorno,
cuja influência no Black--Litterman foi regulada pela
cobertura de notícias e pela confiança estimada. Na
configuração avaliada, o mínimo de três manchetes elegíveis
habilitou a opinião híbrida.

A avaliação financeira considerou carteiras de uma a vinte
ações em dois cenários: grupos de ações comuns aos modelos
e seleção própria por estratégia. Foram realizadas 1.142
simulações na comparação entre Black--Litterman e Markowitz,
além de 1.142 simulações com as referências de prior sem
views e XGBoost direto. As operações incorporaram custos
de negociação e execução no fechamento seguinte ao sinal.

\section{Principais Conclusões e Contribuições}
\label{sec:conclusao_contribuicoes}

Com vinte ativos, o Black--Litterman híbrido apresentou
retorno líquido de 27,79\%, frente a 19,10\% de Markowitz
e 14,87\% da referência livre de risco no período.
O Sharpe foi de 0,70 para o híbrido e 0,29 para Markowitz,
enquanto a queda máxima do patrimônio foi de 8,46\%
e 12,78\%, respectivamente. Nessa configuração, o híbrido
apresentou maior retorno e menor risco realizado que
Markowitz.

Entretanto, o prior sem views obteve retorno de 28,21\%,
Sharpe de 0,74 e menor custo de negociação. Portanto,
a inclusão das opiniões quantitativas e textuais não
melhorou esses indicadores em relação ao prior na
carteira de vinte ativos. Esse resultado evidencia a
importância das expectativas de equilíbrio na configuração
analisada.

Nos grupos de ações comuns, o retorno mediano do
Black--Litterman superou o de Markowitz para todas as
quantidades de ativos a partir de três. Com uma ação,
os resultados coincidiram e com duas, o teto de 50\%
por ativo impôs pesos iguais. A comparação mostrou,
assim, que os limites de concentração também influenciam
as diferenças entre os modelos.

Na seleção própria, o Black--Litterman apresentou maior
retorno em dezenove das vinte quantidades de ações
avaliadas. A exceção foi a carteira com uma única ação,
que perdeu 10,37\%. A separação entre os dois cenários
permitiu distinguir o desempenho da distribuição dos
pesos daquele associado à escolha dos ativos.

A avaliação intermediária mostrou limitações nas previsões
quantitativas. A acurácia direcional média entre as ações
foi de 56,52\%, mas dezenove dos vinte ativos apresentaram
$R^2$ negativo. Esses resultados indicam que o acerto
de direção não foi acompanhado de precisão consistente
na estimativa da magnitude dos retornos.

Na classificação de sentimento, a concordância com a
referência humana definida com apoio de IA foi de 73\%.
A principal dificuldade foi reconhecer manchetes neutras,
cuja revocação alcançou 41,03\%. A comparação identificou
diferenças nos critérios de atribuição de polaridade,
especialmente positiva.

Na avaliação preditiva complementar, a inclusão da
polaridade textual aumentou o erro quadrático de Q,
mas reduziu o erro do posterior BL. A atualização pelo
prior e pela covariância modificou, portanto, o efeito
do sinal textual sobre as expectativas. Já na avaliação
do pregão seguinte ao sinal, as previsões apresentaram
capacidade limitada de antecipar a direção dos retornos.

A cobertura de notícias também restringiu a participação
direta das opiniões: apenas seis dos 240 pares ativo-data
atingiram o mínimo de três manchetes, distribuídos em
três rebalanceamentos mensais. Esse resultado ajuda a
explicar a relevância do prior e destaca a relação entre
a janela de coleta das notícias e a frequência de decisão.

A principal contribuição do trabalho é a implementação
de um fluxo verificável que conecta previsão quantitativa,
classificação textual, controle de confiança e alocação.
A avaliação por etapas permitiu identificar que a vantagem
financeira sobre Markowitz coexistiu com limitações
preditivas e com desempenho próximo ao prior sem views.
A comparação entre grupos comuns e seleção própria
também permitiu examinar separadamente os efeitos da
composição e da distribuição dos pesos.

\section{Trabalhos Futuros}
\label{sec:conclusao_trabalhos_futuros}

Trabalhos futuros podem avaliar os resultados em outros
períodos e condições de mercado, utilizando validação
temporal para selecionar horizontes, atributos, parâmetros
de confiança e quantidade de ações.

A simulação de carteiras utilizou o XGBoost como único
estimador quantitativo de retornos. Em trabalhos futuros,
outros estimadores poderão ser avaliados com os mesmos
ativos, horizontes e períodos, mantendo fixas as regras
de alocação. Isso permitirá verificar se a escolha do
estimador afeta a qualidade das previsões e o desempenho
das carteiras.

Outra extensão é ampliar as comparações entre versões
do Black--Litterman com prior, previsões quantitativas,
sentimento e combinação completa, mantendo as mesmas
condições de execução para examinar a contribuição de
cada componente.

No componente textual, podem ser utilizadas janelas de
notícias mais amplas, ponderadas pelo tempo desde a publicação,
e o conteúdo integral das reportagens. A classificação
também pode ser avaliada em novas amostras por diferentes
avaliadores humanos, com critérios previamente definidos.

A confiança das opiniões pode ser calibrada pelo erro
das previsões, pela qualidade da classificação textual
e pela cobertura de notícias. Também podem ser comparados
rebalanceamentos diários, semanais e mensais.

Na dimensão operacional, o \emph{paper trading} permitiria
acompanhar o sistema em atualização contínua. A simulação
pode incorporar spread, impacto de mercado, tributação
e liquidação das posições, além de dados históricos
de composição do universo e de capitalização.